% Options for packages loaded elsewhere
\PassOptionsToPackage{unicode}{hyperref}
\PassOptionsToPackage{hyphens}{url}
%
\documentclass[
]{article}
\usepackage{amsmath,amssymb}
\usepackage{lmodern}
\usepackage{iftex}
\ifPDFTeX
  \usepackage[T1]{fontenc}
  \usepackage[utf8]{inputenc}
  \usepackage{textcomp} % provide euro and other symbols
\else % if luatex or xetex
  \usepackage{unicode-math}
  \defaultfontfeatures{Scale=MatchLowercase}
  \defaultfontfeatures[\rmfamily]{Ligatures=TeX,Scale=1}
\fi
% Use upquote if available, for straight quotes in verbatim environments
\IfFileExists{upquote.sty}{\usepackage{upquote}}{}
\IfFileExists{microtype.sty}{% use microtype if available
  \usepackage[]{microtype}
  \UseMicrotypeSet[protrusion]{basicmath} % disable protrusion for tt fonts
}{}
\makeatletter
\@ifundefined{KOMAClassName}{% if non-KOMA class
  \IfFileExists{parskip.sty}{%
    \usepackage{parskip}
  }{% else
    \setlength{\parindent}{0pt}
    \setlength{\parskip}{6pt plus 2pt minus 1pt}}
}{% if KOMA class
  \KOMAoptions{parskip=half}}
\makeatother
\usepackage{xcolor}
\setlength{\emergencystretch}{3em} % prevent overfull lines
\providecommand{\tightlist}{%
  \setlength{\itemsep}{0pt}\setlength{\parskip}{0pt}}
\setcounter{secnumdepth}{-\maxdimen} % remove section numbering
\ifLuaTeX
  \usepackage{selnolig}  % disable illegal ligatures
\fi
\IfFileExists{bookmark.sty}{\usepackage{bookmark}}{\usepackage{hyperref}}
\IfFileExists{xurl.sty}{\usepackage{xurl}}{} % add URL line breaks if available
\urlstyle{same} % disable monospaced font for URLs
\hypersetup{
  hidelinks,
  pdfcreator={LaTeX via pandoc}}

\author{Dietmar Gerald Schrausser}
\date{09.05.2026}


\begin{document}

\hypertarget{foliumblat-vr}{%
\section{Foliu{[}m{]}/Blat Vr}\label{foliumblat-vr}}

A \emph{contemporary} English translation is presented from comparing
the existing Latin, Early New High German (ENHG) and English
(translation) texts for better comprehensibility. Especially since the
ENHG text is difficult to understand compared to modern High German, as
a \emph{profound} knowledge of German (as a mother tongue) as well as
already aknowledge of various \emph{dialects} is required.

\hypertarget{on-the-work-of-the-sixth-day}{%
\subsection{On the Work of the Sixth
Day}\label{on-the-work-of-the-sixth-day}}

\begin{quote}
\textbf{S}Exte die dixit de{[}us{]}. p{[}ro{]}ducat terra
a{[}n{]}i{[}m{]}am viue{[}n{]}te{[}m{]}. iume{[}n{]}ta. {[}et{]}
reptila. {[}et{]} bestias terre iuxta spe{[}cie{]}s suas. Et vidit
de{[}us{]} q{[}uo{]}d e{[}ss{]}et bonu{[}m{]}: ait. Faciam{[}us{]}
ho{[}m{]}i{[}n{]}em ad imagine{[}m{]} {[}et{]} si{[}mi{]}litudine{[}m{]}
n{[}ost{]}ram. Et p{[}er{]}sit piscib{[}us{]} mar{[}is{]}. {[}et{]}
volatilib{[}us{]} celi. {[}et{]} bestijs vniuerse terre. Et creuit
de{[}us{]} ho{[}m{]}i{[}n{]}em ad imagine{[}m{]} {[}et{]}
si{[}mi{]}litudine{[}m{]} sua{[}m{]}.
\end{quote}

After the upper parts of the world had been decorated, on the sixth day
he finally decorated the earth with the races of living
beings\footnote{`generib{[}us{]} a{[}n{]}i{[}m{]}aliu{[}m{]}' and
  `geschlechte{[}n{]} der thier'.}. Among the \emph{animals} of the
earth, \textbf{Moses} mentions three: (i) \emph{beasts of burden}, (ii)
\emph{creeping animals}, and (iii) \emph{wild animals}\footnote{`iume{[}n{]}ta.
  reptilia. {[}et{]} bestias' and `iohthier kriechende vnd wildthier'.}.
He thus refers generally to three different types of
\emph{irrational}\footnote{`irr{[}ati{]}onabiliu{[}m{]}' and
  `unvernüftigen'.} animals. \emph{Wild animals} (iii) with perfect
\emph{imagination}\footnote{`pha{[}n{]}tasia' and `fantasey', in
  medieval philosophy, this refers to the imagination.} occupy an
intermediate position among irrational \emph{things}\footnote{`sortite'
  and `thieren'.}, although they can neither be taught nor tamed by
humans. Thus, \emph{crawling animals} (ii) with \emph{im}perfect
\emph{imagination} stand between cattle\footnote{`bruta' and `vieh'.}
and plants. \emph{Beasts of burden} (i) lack reason, but can
nevertheless be largely tamed by humans, so they seem to share in
\emph{some} reason. These occupy, as it were, an intermediate position
between cattle\footnote{`bruta' and `vihe'.} and humans.

Furthermore, he decreed that \emph{living beings}\footnote{`A{[}n{]}i{[}m{]}alia'
  and `thier'.} of every kind and shape, great and small, should be
created in pairs, of both sexes\footnote{`diuersi sex{[}us{]}
  sing{[}u{]}la' and `bede me{[}n{]}dlein vnd freülein'.}. Their
offspring filled the air, earth, and seas. God gave all of them
sustenance from the earth from generation to generation, so that they
could also be \emph{useful} to humans: surplus for food\footnote{`nimi{[}orum{]}
  ad cibos' and `als etlich zu speysung', probably animal
  \emph{products} such as eggs or milk.}, certainly also for
clothing\footnote{`v{[}er{]}o ad vestime{[}n{]}tu{[}m{]}' and `ettlich
  zebeklaidung', like wool, etc.}, those who possess great strength to
help in plowing the land; for this reason, they are also called (in
Latin) \emph{iumenta}\footnote{`vt i{[}n{]} excole{[}n{]}da terra
  iuuarent. vn{[}de{]} dicta sunt iume{[}n{]}ta.'}.\footnote{From the
  `Etymologies' (De Seville,
  \href{https://books.google.com/books?id=wb8Z_vlSyJwC}{1802}, Book XII,
  1.7).}

So far we have spoken of \emph{three} worlds: the (i)
\emph{superheavenly}, the (ii) \emph{celestial}, and the (iii)
\emph{sublunar}. Now we will consider (iv) \emph{humankind} as part of a
fourth world. Having arranged everything in a wondrous way, he decided
to prepare an eternal kingdom for himself, creating countless souls in
advance to whom he could grant immortality. He then created a likeness
of himself, sentient and intelligent, formed in his own
image\footnote{`imaginis sue forma{[}m{]}' and `zu form oder gestalt
  seiner pildnus'.}, to which nothing more perfectly resembles. Thus he
formed humankind from the clay of the earth; hence the name
\emph{human}\footnote{`ho{[}mo{]}' and `mensche{[}n{]}'}, created from
\emph{earth}\footnote{`humo' and `lette{[}n{]} oder kloße der
  erde{[}n{]}'.}.\footnote{reflects a common medieval etymological idea
  that associates \emph{Homo} with \emph{Humus}.}

God, the creator of all things, created man, (which \textbf{Cicero} also
recognized, although he possessed no heavenly scriptures) \emph{this}
man transmitted in the first book on the laws the same as the prophets,
his words we have reproduced below:\footnote{c.f. the `Divinae
  Institutiones' (Lactantius,
  \href{https://books.google.com/books?id=qGsnghsi3uUC}{1497}).}

\begin{quote}
Hoc a{[}n{]}i{[}m{]}al p{[}ro{]}uidu{[}m{]}: sagax: m{[}u{]}ltiplex:
acutu{[}m{]}: memor: plenu{[}m{]} r{[}ati{]}o{[}n{]}is {[}et{]}
{[}con{]}silij. que{[}m{]} vocam{[}us{]} homine{[}m{]}: p{[}re{]}clara
q{[}ua{]}da{[}m{]} {[}con{]}dit{[}i{]}o{[}n{]}e generatu{[}m{]}
e{[}ss{]}e a summo deo solu{[}m{]}.\\
Est e{[}ni{]}m ex tot a{[}n{]}i{[}m{]}antiu{[}m{]} generib{[}us{]}
at{[}que{]} natur{[}is{]}: p{[}ar{]}ticeps r{[}ati{]}o{[}n{]}is {[}et{]}
cogitat{[}i{]}o{[}n{]}is. cu{[}m{]} cetera sint o{[}mn{]}ia
exp{[}er{]}tia.\footnote{`This living being(!) that we call human:
  cautious, agile, versatile, sharp in memory, full of reason and
  counsel; in its clear nature and attributes, born solely of the
  highest God. For it shares in all the races and natures of animate
  creatures, {[}but especially{]} in reason and understanding, which all
  other creatures lack.', `De Legibus' (Cicero \& Cicero,
  \href{https://doi.org/10.3931/e-rara-89116}{1496}, Book I, Section
  22), discusses man as uniquely endowed with reason, which connects him
  with the \emph{divine}.}
\end{quote}

Incidentally, it is a widespread custom among kings or princes of the
earth that, had they conceived and built a \emph{magnificent},
\emph{noble} city, they would, upon its completion, erect their own
image in the city center so that it could be seen and contemplated by
all. Likewise, God, the \emph{Prince} of all things, acted when, after
the entire structure\footnote{`machi{[}n{]}a' and `paw'.} of the world
had been constructed, he finally placed man \emph{in the midst} of it
all, created in his image and formed in his likeness, so that one might
exclaim like \emph{this} \textbf{Mercurius}:\footnote{From `De Vita
  Coelitus Comparanda' (Ficinus,
  \href{https://doi.org/10.3931/e-rara-68641}{1497}).}

\begin{quote}
Magnu{[}m{]} o asclepi mirac{[}u{]}l{[}u{]}m e{[}st{]} ho{[}mo{]}.\\
Hoc p{[}re{]}cipuo no{[}m{]}i{[}n{]}e gl{[}or{]}iari hu{[}m{]}ana
{[}con{]}ditio p{[}otes{]}t: q{[}uo{]} et{[}iam{]}
f{[}a{]}ct{[}u{]}m{]}: vt s{[}er{]}uire illi: nulla creata
s{[}u{]}b{[}stanti{]}a dedignet{[}ur{]}: huic terra et elem{[}en{]}ta:
huic bruta sunt p{[}re{]}sto: {[}et{]} famulant{[}ur{]}. huic militat
celu{[}m{]}: huic salute{[}m{]}: p{[}ro{]}cura{[}n{]}t a{[}n{]}gelices
me{[}n{]}tes.\\
Nec miru{[}m{]} alicui videri debet: amari illu{[}m{]} ab
om{[}n{]}ibus.\\
In quo om{[}n{]}ia suu{[}m{]} aliq{[}ui{]}d:\\
Immo se tota {[}et{]} sua om{[}n{]}ia agnoscu{[}n{]}t.\footnote{`O
  Asclepius, man is a great wonder. Through this extraordinary name,
  human existence can boast: for it has come to pass that no creature
  disdains to serve him. Earth and elements stand by him; the animals
  are ready for him and serve him. For him the heavens fight; for him
  angels bring salvation. It should surprise no one that he is loved by
  all, in whom everything finds something of its own; indeed, everything
  and everything of its own recognizes in him', from the `Asklepios', a
  Hermetic text (vgl. Trismegistu,
  \href{https://gdz.sub.uni-goettingen.de/id/PPN858183994}{1471}), also
  quoted in in Pico's `Oratio de hominis dignitate' (Mirandola \&
  Mirandola, \href{https://doi.org/10.3931/e-rara-61217}{1601}).}
\end{quote}

\hypertarget{references}{%
\subsection{References}\label{references}}

Cicero, M. T., \& Cicero, Q. (1496). \emph{M.T. Ciceronis de Natura
Deorum Libri Tres} \ldots{} Impræssum Venetiis: per Symonem Papiensem
dictum Biuilaqua. \url{https://doi.org/10.3931/e-rara-89116}

De Seville, I. (1802). \emph{S. Isidori Hispalensis Episcopi Hispaniarum
Doctoris Opera Omnia}. Edited by Arevalo, F. Romae: Apvd Antonivm
Fvlgonivm. \url{https://books.google.com/books?id=wb8Z_vlSyJwC}

Ficinus, M. (1497). \emph{De Triplici Vita}. Basel: Johann Amerbach.
\url{https://doi.org/10.3931/e-rara-68641}

Lactantius, L.C.F. (1497). \emph{Divinae Institutiones}. Venetiis: Simon
Bevilaqua. \url{https://books.google.com/books?id=qGsnghsi3uUC}

Mirandola, G., \& Mirandola, G. F. (1601). \emph{Ioannis Pici,
Mirandulae \ldots{} opera quae extant omnia \ldots{}} Basileae: per
Sebastianum Henricpetri. \url{https://doi.org/10.3931/e-rara-61217}

Trismegistu, H. (1471). \emph{Mercurii Trismegisti Liber de Potestate Et
Sapientia Dei e Graeco in Latinum Traductus a Marsilio Ficino Pimander
Incipit}. Edited by Ficinus, M. Treviso.
\url{https://gdz.sub.uni-goettingen.de/id/PPN858183994}

\end{document}
